\documentclass[11pt]{article}

%================================================================
% Paquetes
%================================================================
\usepackage[utf8]{inputenc}
\usepackage[T1]{fontenc}
\usepackage[spanish,shorthands=off]{babel} % shorthands=off evita que "<" y ">"
                                            % queden activos y choquen con las
                                            % flechas de TikZ (->, <->, etc.)
\usepackage{pgfplots}
\usepackage[margin=2.5cm]{geometry}
\usepackage{amsmath,amssymb}
\usepackage{mathpazo}
\usepackage{xcolor}
\usepackage{tcolorbox}
\tcbuselibrary{skins,breakable}
\usepackage{tikz}
\usepackage{enumitem}
\usepackage{titlesec}
\usepackage{caption}
\usepackage{multirow}
\usepackage{graphicx}
\usepackage{booktabs}
\usepackage{bm}
%================================================================
% Paleta de colores
%================================================================
\definecolor{ucred}{RGB}{117,26,41}    % granate, tono institucional
\definecolor{devcolor}{RGB}{45,70,95}  % azul pizarra, para el desarrollo


\usepackage{listings}

\definecolor{codebg}{RGB}{248,248,248}
\definecolor{codekw}{RGB}{0,92,197}      % palabras clave (def, for, if...)
\definecolor{codestr}{RGB}{163,21,21}    % strings
\definecolor{codecom}{RGB}{106,153,85}   % comentarios
\definecolor{codenum}{RGB}{9,134,88}     % números

\lstdefinestyle{pythonstyle}{
    backgroundcolor=\color{codebg},
    basicstyle=\ttfamily\small,
    keywordstyle=\color{codekw}\bfseries,
    stringstyle=\color{codestr},
    commentstyle=\color{codecom}\itshape,
    numberstyle=\tiny\color{gray},
    numbers=left,
    numbersep=8pt,
    frame=single,
    rulecolor=\color{gray!40},
    breaklines=true,
    breakatwhitespace=false,
    showstringspaces=false,
    tabsize=4,
    language=Python,
    morekeywords={self},
    xleftmargin=1.5em,
    framexleftmargin=1.2em,
    inputencoding=utf8,          % <- nuevo
    extendedchars=true,          % <- nuevo
    literate=%                   % <- nuevo: le dice a listings cómo dibujar cada tilde
        {á}{{\'a}}1 {é}{{\'e}}1 {í}{{\'i}}1 {ó}{{\'o}}1 {ú}{{\'u}}1
        {Á}{{\'A}}1 {É}{{\'E}}1 {Í}{{\'I}}1 {Ó}{{\'O}}1 {Ú}{{\'U}}1
        {ñ}{{\~n}}1 {Ñ}{{\~N}}1 {ü}{{\"u}}1 {Ü}{{\"U}}1
        {¿}{{?`}}1 {¡}{{!`}}1 {·}{{$\cdot$}}1
        {°}{{$^\circ$}}1 {µ}{{$\mu$}}1 {×}{{$\times$}}1 {⊥}{{$\perp$}}1
        {-}{{-\allowbreak}}1,  % deja partir las lineas largas de guiones
}

\setcounter{MaxMatrixCols}{12} 
%================================================================
% Estilo de secciones
%================================================================
\titleformat{\section}
  {\normalfont\Large\bfseries\color{ucred}}{\thesection}{0.6em}{}
\titlespacing*{\section}{0pt}{1.4em}{0.8em}
\usetikzlibrary{patterns, arrows.meta}


\usetikzlibrary{shapes.geometric}


%================================================================
% Cajas (tcolorbox)
%================================================================

% Caja para el enunciado del problema
\newtcolorbox{enunciado}[1]{
  enhanced, breakable,
  colback=ucred!4, colframe=ucred,
  coltitle=white, fonttitle=\bfseries\large,
  title={#1},
  boxrule=0.9pt, arc=2mm,
  top=8pt, bottom=10pt, left=9pt, right=9pt
}

% Caja para que el estudiante desarrolle su solucion
\newtcolorbox{desarrollo}[1]{
  enhanced, breakable,
  colback=white, colframe=devcolor,
  coltitle=white, fonttitle=\bfseries,
  title={#1},
  boxrule=0.8pt, arc=1.5mm,
  top=10pt, bottom=10pt, left=9pt, right=9pt
}

%================================================================
% Espacios para completar (Tarea 3)
%================================================================
\definecolor{pendcolor}{RGB}{190,120,20}  % ambar, para lo que falta

% Caja de lo que falta escribir (analisis, discusion, valores).
% Uso: \begin{pendiente} ... \end{pendiente}   -> borrarla cuando se complete
\newtcolorbox{pendiente}[1][Por completar]{
  enhanced, breakable,
  colback=pendcolor!6, colframe=pendcolor,
  coltitle=white, fonttitle=\bfseries\small, fontupper=\small,
  title={#1}, boxrule=0.6pt, arc=1mm, borderline={0.4pt}{0pt}{pendcolor,dashed},
  top=5pt, bottom=5pt, left=7pt, right=7pt
}

% Figura que se completa sola: si el archivo existe lo inserta, si no deja un
% recuadro con lo que va ahi y el nombre de archivo esperado.
% Uso: \figuraT{ruta/archivo.png}{ancho (fraccion de \linewidth)}{que debe mostrar}
\newcommand{\figuraT}[3]{%
  \IfFileExists{#1}{%
    \includegraphics[width=#2\linewidth]{#1}%
  }{%
    \begin{tikzpicture}
      \draw[pendcolor, dashed, thick, rounded corners=2pt, fill=pendcolor!4]
        (0,0) rectangle (#2\linewidth, 5.5cm);
      \node[text width=#2\linewidth-1cm, align=center, font=\small, text=pendcolor!80!black]
        at (#2\linewidth/2, 2.75cm)
        {\textbf{Figura pendiente}\\[3pt] #3\\[6pt]
         {\scriptsize archivo esperado: \texttt{\detokenize{#1}}}};
    \end{tikzpicture}%
  }%
}

% Celda de tabla por completar
\newcommand{\pend}{\textcolor{pendcolor}{\textbf{?}}}

\newcommand{\dnode}[5]{\coordinate (#1) at ({#2 + (\amp*#4)/1000},{#3 + (\amp*#5)/1000});}

% ===== MACRO: color proporcional al esfuerzo =====
% Escala roja (tracción) hasta el máximo, azul (compresión) hasta el mínimo
% Uso: \sigmabar{nodoA}{nodoB}{valor_sigma}
\newcommand{\sigmabar}[3]{%
  \pgfmathsetmacro{\sval}{#3}%
  \pgfmathsetmacro{\aval}{abs(\sval)}%
  \ifdim \aval pt < 1pt
    \draw[line width=1.4pt, gray] (#1) -- (#2);
  \else
    \ifdim \sval pt < 0pt
      \pgfmathsetmacro{\pct}{100*\aval/\smin}
      \draw[line width=1.4pt, blue!\pct!white] (#1) -- (#2);
    \else
      \pgfmathsetmacro{\pct}{100*\aval/\smax}
      \draw[line width=1.4pt, red!\pct!white] (#1) -- (#2);
    \fi
  \fi
}

%================================================================
% Documento
%================================================================
\begin{document}

%---------------------- Encabezado -----------------------------
\begin{center}
    {\scshape\large Pontificia Universidad Cat\'olica de Chile}\\[2pt]
    {\scshape Escuela de Ingenier\'ia}\\[2pt]
    {\small Departamento de Ingenier\'ia Mec\'anica y Metal\'urgica}\\[2pt]
    {\small ICM 3323 --- Elementos Finitos en Dise\~no Mec\'anico}\\[10pt]
    {\color{ucred}\rule{0.42\linewidth}{0.6pt}}\\[10pt]
    {\Huge\bfseries\color{ucred} TAREA N\textdegree\ 3}
\end{center}

\vspace{4pt}
\begin{flushright}
\renewcommand{\arraystretch}{1.4}
\begin{tabular}{rl}
    Nombre: & Pedro Pablo Taboada Santiba\~nez \\
    Fecha:  & \today\\
\end{tabular}
\end{flushright}

\vspace{0.6cm}

%---------------------- Problema ---------------------------------
\section*{Problema}

\begin{enunciado}{Problema (6 puntos)}
    \begin{enumerate}[label=\alph*), itemsep=4pt]
        \item Calcular mediante el m\'etodo de elementos finitos los campos de
        las diferentes variables correspondientes al problema del apartado 35
        del libro \textit{``Theory of Elasticity''} de S.P.~Timoshenko y
        J.N.~Goodier (todos los par\'ametros son de libre elecci\'on):
        \textit{``The Effect of Circular Holes on Stress Distributions in
        Plates''} (figura 52).
        \item Comparar y discutir los resultados obtenidos en a) con los
        correspondientes a la soluci\'on anal\'itica del mencionado apartado.
        \item Comparar y discutir los resultados obtenidos en a) con el
        proporcionado por un programa comercial de elementos finitos.
    \end{enumerate}
\end{enunciado}

\begin{figure}[!htbp]
\centering
\begin{tikzpicture}[>=latex, scale=1.1]
    % --- Datos geometricos (solo para el dibujo) ---
    \def\Lplaca{9}      % largo de la placa
    \def\dplaca{3}      % ancho de la placa (d)
    \def\rag{0.6}    % radio del agujero (r)
    \pgfmathsetmacro{\xc}{\Lplaca/2}
    \pgfmathsetmacro{\yc}{\dplaca/2}

    % placa con agujero
    \draw[thick, fill=black!5, even odd rule]
        (0,0) rectangle (\Lplaca,\dplaca)
        (\xc,\yc) circle (\rag);

    % ejes de simetria
    \draw[dashed, thin, gray] (-0.2,\yc) -- (\Lplaca+0.2,\yc);
    \draw[dashed, thin, gray] (\xc,-0.3) -- (\xc,\dplaca+0.3);

    % traccion uniforme S en ambos extremos
    \foreach \y in {0.25,0.75,...,2.75}{
        \draw[-{Latex[length=2mm]}, thick] (0,\y) -- (-0.8,\y);
        \draw[-{Latex[length=2mm]}, thick] (\Lplaca,\y) -- (\Lplaca+0.8,\y);
    }
    \node at (-1.2,\yc) {\Large $S$};
    \node at (\Lplaca+1.2,\yc) {\Large $S$};

    % puntos m (borde de la placa) y n (borde del agujero)
    \fill (\xc,\dplaca) circle (1.6pt) node[above=3pt] {$m$};
    \fill (\xc,\yc+\rag) circle (1.6pt) node[above left=1pt] {$n$};

    % cota 2r
    \draw[<->, thin] (\xc+1.1,\yc-\rag) -- (\xc+1.1,\yc+\rag)
        node[midway, fill=black!5, inner sep=2pt] {$2r$};
    \draw[thin] (\xc,\yc-\rag) -- (\xc+1.25,\yc-\rag);
    \draw[thin] (\xc,\yc+\rag) -- (\xc+1.25,\yc+\rag);

    % cota d
    \draw[<->, thick] (\Lplaca+2.0,0) -- (\Lplaca+2.0,\dplaca)
        node[midway, fill=white, inner sep=2pt] {$d$};
    \draw[thin] (\Lplaca+0.1,0) -- (\Lplaca+2.2,0);
    \draw[thin] (\Lplaca+0.1,\dplaca) -- (\Lplaca+2.2,\dplaca);
\end{tikzpicture}
\caption{Placa de ancho $d$ con un agujero circular de di\'ametro $2r$ en su
eje de simetr\'ia, sometida a tracci\'on uniforme $S$ en sus extremos
(Timoshenko \& Goodier, Fig.~52).}
\label{fig:t3-placa}
\end{figure}

%---------------------- Desarrollo -------------------------------

\begin{desarrollo}{a) Soluci\'on por MEF de los campos de la placa con agujero}
    % Par\'ametros elegidos (S, r, d, largo, espesor, E, nu), modelo (tensi\'on
    % plana / deformaci\'on plana, simetr\'ia), elemento, malla, condiciones de
    % borde, c\'odigo y resultados (desplazamientos, sigma_x, sigma_y, tau_xy,
    % sigma_r, sigma_theta, tau_rtheta).

\end{desarrollo}

\begin{desarrollo}{b) Comparaci\'on con la soluci\'on anal\'itica}
    % Soluci\'on de Kirsch (Ec. 61), sigma_theta en el borde del agujero,
    % sigma_theta en theta = pi/2 vs r, factor de concentraci\'on en n y m,
    % efecto del ancho finito d.

\end{desarrollo}

\begin{desarrollo}{c) Comparaci\'on con un programa comercial de elementos finitos}
    % Programa usado (ANSYS, Abaqus, SolidWorks, ...), modelo, malla,
    % resultados y comparaci\'on con a) y b).

\end{desarrollo}


\end{document}
